\documentclass[10pt]{article}
\usepackage[margin=0.95in]{geometry}
\usepackage{amsmath,amssymb,booktabs,multicol}
\usepackage[T1]{fontenc}
\usepackage{microtype}
\setlength{\parskip}{2pt}
\title{\vspace{-1.3em}\large\textbf{Specification note: the E-series architecture --- equivalence-scale preferences, sequential fertility, and honest income risk}\vspace{-0.4em}}
\author{\normalsize Fertility--housing project}
\date{\normalsize July 20, 2026}
\begin{document}
\maketitle
\vspace{-1.5em}

\paragraph{The change in one view.}
The E-series replaces the production model's two deepest provisional
structures at once. Preferences lose their Stone--Geary intercepts ---
children price in through an expenditure-share shifter and a
utility-level scale instead of committed bundles --- and fertility becomes
sequential: households try for a first child and one-child households may
later try for a second, each attempt succeeding with an age-declining
fecundity probability. Removing the intercepts removes the feasibility
bind that capped income risk, so the income process is set externally at
its literature value rather than tuned down to keep budget sets nonempty.

\begin{center}\footnotesize
\begin{tabular}{lll}
\toprule
 & Production (M5) & E-series \\
\midrule
Child costs & floors $\bar c(n), \bar h(n)$ (5 params, SMM) & $\alpha(n,s)$, $\tilde e(n,s)$ (3 params, SMM) \\
Feasibility & bundle must be affordable; gate + censoring & every $y>0$ feasible by construction \\
Income risk & $\sigma_z=0.0645$ (feasibility-capped) & $\sigma_z=0.20$, $\rho_z=0.960$ (external) \\
Fertility & one-shot family-size choice & sequential attempts, fecundity $\pi_j$ (external) \\
Free params / moments & 14 / 15 & 12 / 15 \\
\bottomrule
\end{tabular}
\end{center}

\paragraph{Preferences.}
Period utility is
\[
u(c,h;n,s)=\tilde e(n,s)\,
\frac{\bigl(c^{\alpha(n,s)}h^{1-\alpha(n,s)}\bigr)^{1-\sigma}}{1-\sigma}
+\psi n_k,\qquad
\tilde e = 1+\gamma_e n_k,\quad
\alpha = \alpha_0-(\delta^{jump}_\alpha+\delta_\alpha n_k)\mathbf{1}\{n_k>0\},
\]
with $\sigma=2$ maintained, so $\tilde e$ coincides with the equivalence
scale over the composite: a family needs $\tilde e$ times the childless
composite for equal welfare. Textbook good-specific (Barten) scales are
deliberately not used: under a Cobb--Douglas composite they cancel from
the consumption--housing allocation, which would erase the child--housing
demand channel the paper is about; the share shifter $\alpha(n,s)$
restores that channel in the form the consumption literature uses
(demographic taste shifters), and $\tilde e$ carries the cost of children.
With no intercepts, marginal utility diverges at zero and any positive
income admits a feasible choice.

How the references handle the same object: Adda, Dustmann and Stevens
(2017, JPE) \emph{impose} the OECD-modified scale (1 + 0.5 per adult +
0.3 per child) on household consumption; Scholz, Seshadri and
Khitatrakun (2006, JPE) impose the Citro--Michael scale
$(A+0.7K)^{0.7}$ in utility \emph{and} in their means-tested transfer
floor; Fern\'andez-Villaverde and Krueger (2007, ReStat) deflate by the
mean of six published scales; Blundell, Pistaferri and Saporta-Eksten
(2016, AER) estimate demographic preference shifters inside their
structural system; Borella, De Nardi and Yang (2023, ReStud) and
De Nardi, Fella and Paz-Pardo (2025, ReStud) use divisor scales
varying with marital status and children --- the former imposing the
Citro--Michael scale $(j+0.7f)^{0.7}$, the latter imposing the
OECD-modified weights (1, 0.5 per second adult, 0.3 per child).
Guner, Kaygusuz and Ventura (2020, ReStud) explicitly
reject utility scaling --- but only to preserve balanced-growth
consistency, a constraint our stationary economy does not face.
Identification discipline for scales is the Pollak--Wales (1979, AER)
critique: welfare-level scales are not identified from demand data alone.
\textbf{Our approach differs from all of the above: we estimate
$\gamma_e$ inside the SMM} from housing, fertility, and wealth moments
--- and the estimate, $\gamma_e = 0.203$, implies $+20.3\%$ per one child
and $+40.5\%$ for two, reproducing the OECD-modified child weights
($+20\%/+40\%$ on a couple base) that the literature imposes. Because our
$\gamma_e$ is identified by behavior through the model structure, this
agreement is corroboration rather than circularity.

\paragraph{Fertility technology.}
The production model chooses completed family size in one event; the
E-series prices time. Childless households choose \{wait, try\}; one-child
households with the window open choose \{stop, try for the second\}; an
attempt succeeds with probability $\pi_j = \max\{0,\min\{1,
1-\omega_1e^{\omega_2(a_j-18)}\}\}$, zero from 45, and a second birth
restarts the single child clock (per-child age tracking stays collapsed
--- the reason sequencing was originally dropped --- at the small cost
that the first child's maturation rides the youngest's clock). The
schedule is \emph{external}: biology identifies it, and pinning it keeps
the fertility logit scale $\kappa_f$ identified, since both objects
smooth hazards. Calibration precedents and our fit are documented in the
companion note (\texttt{fecundity\_schedule\_note\_20260720}): Sommer
(2016, JME) fits an exponential permanent-infertility CDF to
Trussell--Wilson (1985), reaching $\approx$97\% infertile at 45; de la
Croix and Pommeret (2021, JET) fit a bounded logistic to L\'eridon's
conception-by-duration table with a 4\% natural-sterility level from
Baudin, de la Croix and Gobbi (2015, AER). Our deployed
$(\omega_1,\omega_2)=(0.02,0.134)$ evaluates to $\pi=\{0.97, 0.90, 0.81,
0.62\}$ at ages 22/30/35/40 against L\'eridon's (2004)
conceive-within-four-years anchors $\{0.96, 0.91, 0.84, 0.64\}$ --- on
the evidence to 1--3 points; the formal least-squares fit is a pending
half-day refinement. New measured objects: attempt and realized birth
hazards by age, first-birth age distribution, birth spacing via the flow
progression from one to two children, and the decomposition of
childlessness into chosen versus ran-out-of-clock.

\paragraph{Income process.}
Production carries $\sigma_z=0.0645$ --- a third of the PSID range ---
because the Stone--Geary bundle had to be affordable in every state (the
July-18 frontier computation: infeasible above $\sigma_z\!\approx\!0.09$
at the calibrated persistence). With no intercepts the constraint is
gone, and the E-series imposes $(\rho_z,\sigma_z)=(0.960, 0.20)$
externally (the Sommer--Sullivan range), with no transfer floor needed.

\paragraph{Calibration accounting and status.}
Twelve free parameters (the five floors leave, $\delta^{jump}_\alpha,
\delta_\alpha, \gamma_e$ enter) on the unchanged 15-moment
income-disciplined target system; externals: $\theta_n=0$, $\sigma=2$,
$\phi=0.80$, the income process, and the fecundity schedule. The share
shifters are identified by the family/childless housing and rooms
moments; $\gamma_e$ by fertility levels jointly with wealth accumulation.
Current status (E2, July 20): loss 5.49 versus the production model's
9.04 on the same moments, cross-evaluator confirmed; old-age ownership
and the estate distribution resolved; the open issues are young renters'
liquid wealth (0.52 vs.\ 0.18) and the pending fecundity fit plus
timing-moment data targets before paper-facing use.

\begin{footnotesize}
\begin{multicols}{2}
\noindent Adda, Dustmann, Stevens (2017), \emph{JPE} 125(2).\\
Baudin, de la Croix, Gobbi (2015), \emph{AER} 105(6).\\
Blundell, Pistaferri, Saporta-Eksten (2016), \emph{AER} 106(2).\\
Borella, De Nardi, Yang (2023), \emph{ReStud} 90(1).\\
de la Croix, Pommeret (2021), \emph{JET} 193.\\
De Nardi, Fella, Paz-Pardo (2025), \emph{ReStud} 92(2).\\
Fern\'andez-Villaverde, Krueger (2007), \emph{ReStat} 89(3).\\
Guner, Kaygusuz, Ventura (2020), \emph{ReStud} 87(5).\\
L\'eridon (2004), \emph{Human Reproduction} 19(7).\\
Pollak, Wales (1979), \emph{AER} 69(2).\\
Scholz, Seshadri, Khitatrakun (2006), \emph{JPE} 114(4).\\
Sommer (2016), \emph{JME} 83.\\
Trussell, Wilson (1985), \emph{Population Studies} 39(2).
\end{multicols}
\end{footnotesize}

\end{document}
